% Build the layer, Chapter 7, first build: a ledger that conserves units.
% Inserted at the end of \S\ref{sec:conservation}, after the property on
% where a forfeit goes. It builds the ledger that the Conservation theorem
% describes. The paragraph "The implementation" says that Port Daddy's
% ledger is not that ledger yet.
%
% No code or output in this file is typed by hand. \pdexcerpt, \pdrun and
% \pdrunlines print files under excerpts/, which the kernel's
% tools/excerpts.mjs cuts from its source and its recorded runs (see
% figures/pd-listing.tex). The figure draws two of those recorded tables as
% printed. To change a listing or the figure, change the kernel and cut the
% excerpts again. The prose refers to functions by name and never by line
% number, so a listing can move within its file.
%
% \pdlistingroom{n} before a listing of n lines keeps it on one page. When a
% listing changes length, change the number to match.
\subsection{Build the layer: a ledger that conserves units}
\label{sec:bonded-build}

\noindent\emph{Express lane: this section builds the ledger of
Theorem~\ref{thm:conservation} twice. The first ledger stores dollars in
floating point, and it refuses a bond that its wallet covers. The second
stores whole units, and every subtraction in it is checked. An exhaustive
test visits every state of it within the bounds of this chapter's
TLA$^{+}$ model, and Kani checks the step of the theorem's induction for
any 64-bit amounts. With one guard removed, units are still conserved and
one bond is refunded twice.}

\S\ref{sec:conservation} states its theorem for money counted in whole
units. It then says that Port Daddy's ledger stores dollars in floating
point, so the theorem describes an intended ledger and not the running
one. This section builds the intended ledger and runs it. The code is
\pdkernelname{}, the small kernel written for this book, which gains a
layer in each chapter. It is not the daemon whose ledger the rest of this
chapter describes.

The section has the same four parts as
\weld{prereq:sec:prereq-build}{Naive}: \emph{Naive}, \emph{Break},
\emph{Fix} and \emph{Verify}. As there, a script cuts each listing from
the kernel's source, and each run is a recording. The layer is small. It
holds one project's funds in memory and decides four operations on them.
It does not authenticate the caller of an operation, and the adjudicator
of \S\ref{sec:governance} is not built.

\subsubsection{Naive: three balances in dollars}

An object that we will call the \emph{ledger} holds a project's funds. It
has the three buckets of Definition~\ref{def:acct-state} and a list of
bonds, and it has the theorem's operations under shorter names.
\texttt{topUp} adds new money to the wallet. \texttt{post} moves the
amount of a bond from the wallet into escrow, which the theorem writes as
$\mathrm{escrow}(b)$. \texttt{refund} and \texttt{slash} settle a bond.
The ledger handles one request at a time and finishes it before it starts
the next.

Prices are quoted in dollars, so the first ledger stores dollars
(Listing~\ref{alg:bonded-build-float}).

\pdlistingroom{14}
\pdexcerpt[caption={The first ledger posts a bond. Its balances are
dollars in floating point.},label={alg:bonded-build-float}]{ch07-float}

This ledger is written with care. It compares the wallet with the bond and
debits the wallet in the same step, so no other request can run between
the two. It records whether a bond was settled, and it will not forfeit
more than a bond holds. The runs below pass every ledger its amounts in
cents, and this one divides them by 100.

\subsubsection{Break: thirty cents}

Alice tops a project up with thirty cents. Wren posts a bond of ten cents
from it, and then a bond of twenty.

\pdlistingroom{14}
\pdrun[caption={The break. The wallet held thirty cents, and the two bonds
add up to thirty.},label={alg:bonded-build-run-float}]{thirty-cents-float}

After the first bond the wallet holds 0.19999999999999998 dollars. That is
less than 0.2, and the second bond is refused. No line of
Listing~\ref{alg:bonded-build-float} is a bug in the usual sense. Its
numbers are binary fractions of 53 bits, and a tenth has no finite
expansion in binary, as a third has none in
decimal~\pdcite{goldberg1991floating}. The ledger stores the nearest
fraction it can, and the error appears after a subtraction. Here the
error is under $3\times10^{-17}$ of a dollar, and it is enough to refuse
the bond. \emph{Now you try:} Exercise~\ref{ex:bonded-build-cents} derives
the refusal by hand, and then finds two bonds that are refused in one
order and posted in the other.

Port Daddy's tests compare the three buckets with the expected supply to
six decimal places, and they pass. They would pass here too. A test of
totals does not detect a comparison that gave the wrong result, and the
\emph{Verify} part below counts how often that happens.

\subsubsection{Break again: two workers and one wallet}

Storing cents fixes that. A second question remains: where the wallet is
compared with the bond. Suppose a ledger stores whole units and has two
requests. One reports whether the wallet holds a given amount, and the
other moves the units. A careful worker checks before it posts
(Listing~\ref{alg:bonded-build-check-then-post}).

\pdlistingroom{9}
\pdexcerpt[caption={A worker that checks the wallet, and then
posts.},label={alg:bonded-build-check-then-post}]{ch07-check-then-post}

The \texttt{yield} is where the worker stops between its two requests and
another worker may run (\weld{prereq:sec:prereq-build}{yield}). Heron and
Wren each want a bond of 6 from a wallet of 10. The wallet covers one bond
and not two. \pdkernelname{} runs every order in which the two workers can
take their turns.

\pdlistingroom{12}
\pdrun[caption={Every order of two workers' turns, on a ledger that leaves
the check to the worker.},label={alg:bonded-build-posts-unchecked}]{posts-unchecked}

Four of the six orders post both bonds, and the wallet ends at $-2$. They
are the four orders in which both workers check before either posts. This
is Chapter~0's lost update, with a wallet in place of the counter
(\weld{prereq:sec:prereq-build}{lose an update}).
Theorem~\ref{lem:no-overdraft} excludes it by assuming that the check and
the debit commit together. The first ledger satisfied that assumption and
stored inexact numbers. This one stores exact numbers and does not satisfy
it.

\subsubsection{Fix: whole units, and a checked subtraction}

The fix is in the kernel's decision core, the Rust function that decides
every proposal in this book's code
(\weld{prereq:sec:prereq-build}{decision core}). Its funds are three counts
(Listing~\ref{alg:bonded-build-funds}).

\pdlistingroom{16}
\pdexcerpt[caption={A project's funds, and the total number of
units.},label={alg:bonded-build-funds}]{ch07-funds}

A \texttt{u64} is a whole number, 0 or greater. It cannot represent a
fraction of a unit or a negative balance. The unit is the project's
smallest amount, which in these runs is a cent. \texttt{supply} is $S_Q$
of Definition~\ref{def:acct-state}. It returns a value only when the
buckets add up to no more than \texttt{CEILING}, which is $2^{53}-1$. The
kernel's host is JavaScript, which stores every number as a double, and
above $2^{53}$ a double cannot represent every whole number. A count that
was exact in the core and rounded in the host would be the first ledger's
fault again, one layer up. So the core refuses a top-up that would exceed
the ceiling.

The wallet has one guard, \texttt{covered}
(Listing~\ref{alg:bonded-build-covered}).

\pdlistingroom{11}
\pdexcerpt[caption={The guard on the wallet. It returns the wallet's new
balance.},label={alg:bonded-build-covered}]{ch07-covered}

\texttt{checked\_sub} subtracts, and returns nothing if the result would
be negative. So \texttt{covered} does not return yes or no. It returns the
wallet's new balance. The line marked \texttt{MUTANT} belongs to a second
build of the core, which has no guard. The \emph{Verify} part runs it.

Posting a bond is Listing~\ref{alg:bonded-build-post}.

\pdlistingroom{22}
\pdexcerpt[caption={The decision to post a
bond.},label={alg:bonded-build-post}]{ch07-post}

\texttt{post} gets the wallet's new balance from \texttt{covered}, and it
has no other way to get that number. No line debits the wallet without
the comparison, because the two are one expression.

A refund and a slash are one function, \texttt{settle}
(Listing~\ref{alg:bonded-build-settle}).
Table~\ref{tab:bonded-conservation-dictionary} treats a refund the same
way, as a slash of zero. The bond's amount leaves escrow, and the wallet
and the pool together receive the same amount, so the total cannot change
here.

\texttt{settle} has three guards. The first is \texttt{held}
(Listing~\ref{alg:bonded-build-held}): the bond exists, and it has not
been settled. The other two work like \texttt{covered}.
\texttt{after\_forfeit} subtracts the forfeit from the bond, so that no
more is forfeited than the bond holds. \texttt{backed} subtracts the bond
from escrow. The kernel keeps escrow equal to the sum of the held bonds,
and the core checks anyway. The core receives its state from the kernel
and cannot verify that the state is right. \texttt{settle} is passed the
forfeit as a whole number of units. Exercise~\ref{ex:bonded-build-split}
asks who gets the remainder when half of a bond of 5 is forfeited.

\pdlistingroom{21}
\pdexcerpt[caption={The decision to settle a bond. A refund is a slash
that forfeits nothing.},label={alg:bonded-build-settle}]{ch07-settle}

\pdlistingroom{13}
\pdexcerpt[caption={The guard on a bond's
status.},label={alg:bonded-build-held}]{ch07-held}

The ledger that holds the funds is short
(Listing~\ref{alg:bonded-build-core-ledger}).

\pdlistingroom{23}
\pdexcerpt[caption={The kernel's ledger. It sends each operation to the
core, and stores the result before it
returns.},label={alg:bonded-build-core-ledger}]{ch07-core-ledger}

It passes the core the funds and the bond that an operation names, and it
stores the core's result. The check and the change are one step for the
same reason as in Chapter~0's cell: no other request runs between
\texttt{decide} and the line that stores its result.

% A page may break here. See the note in anchor-build-the-verifier.tex.
\pagebreak[0]
\keyidea[Make the check the only way to subtract]{A balance is safe when
no code can lower it except through the check. A check that returns yes or
no leaves the subtraction to a later line, and that line can be written or
changed without the check. A check that returns the new balance cannot be
skipped.}

\Needspace{12\baselineskip}
\subsubsection{Verify: every state, and then any amount}

First, the thirty cents.

\pdlistingroom{9}
\pdrun[caption={The fix. Both bonds are posted, and the wallet is
empty.},label={alg:bonded-build-run-core}]{thirty-cents-core}

Next, the settlement that \S\ref{sec:conservation} works by hand.

\pdlistingroom{10}
\pdrun[caption={The chapter's worked settlement, run. Of a bond of 6, 2 is
forfeited.},label={alg:bonded-build-settlement}]{settlement-core}

Step 2 leaves the project at $(4,6,0)$ and step 3 at $(8,0,2)$, as in the
example. Then the two workers, on every ledger.

\pdlistingroom{8}
\pdrun[caption={Two workers and one wallet, on each
ledger.},label={alg:bonded-build-posts}]{posts}

\paragraph{Every state.}
Three runs are not much evidence. TLC checked this chapter's model in
every state it can reach with amounts of 1 to 3 and at most 6 units
minted. \pdkernelname{} tests its own ledger the same way, with those
amounts and at most 3 bonds posted, where the model has 3 agents. The test
is exhaustive within those bounds: in each state it applies every
operation the bounds allow, and it continues from each new state. Some of
those operations should be refused: a bond the wallet cannot cover, or a
forfeit of one unit more than a bond holds. The test also tries to settle
every bond that was ever posted, whether it is held or not, and one that
never existed. Each result is compared with a reference model, which is
the proof sketch of Theorem~\ref{thm:conservation} written as arithmetic
on whole numbers in a file that shares no code with the core.

\pdlistingroom{17}
\pdrun[caption={Every state each ledger can reach, compared with the
reference model.},label={alg:bonded-build-states}]{states}

The kernel's ledger reaches 6,902 states. In every one of them the buckets
hold exactly what was topped up. It accepts all 25,672 operations that the
model allows, and it refuses every operation that the model refuses. The
ledger that leaves the check to the worker reaches the same states, and
accepts 653 bonds that its wallet cannot cover.

The ledger in floating point reaches more states, because the same funds
are stored as different doubles after different sequences of operations.
It refuses 386 bonds that its wallet covers. In 31,174 of its 36,662
states a bucket differs from the model's, and in 18,873 the buckets do not
add up to the top-ups. The largest of those differences is
$2.8\times10^{-17}$ of a dollar. A test to six decimal places passes in
every one of these states.

Figure~\ref{fig:bonded-moves} draws the operations that the kernel's
ledger accepted in the states with all 6 units topped up.

% Brief (grammar: every state as a point, every move as a stroke; atlas
%   VI/fig:bonded-moves).
% Reader question: from where a project's funds are, where can a move take
% them?
% Claim: no move the kernel makes leaves the triangle of funds that add up
% to what was topped up. A post or a refund moves along a row, a slash moves
% up a row, and nothing moves down. With the check of the wallet removed,
% moves leave the triangle.
% Evidence: every move made from a state with all 6 units topped up, in the
% walk of every state within the bounds [internal, harbor states --each;
% sessions states-each.txt and states-each-overdraw.txt, drawn as printed,
% and tests/harbor-research/test_bonded_ledger_numbers.py holds the counts
% in the caption to them].
% Must distinguish: a move along a row from a move up; a move the Book's
% rules allow from one they refuse; inside the triangle from outside it.
% Grammar chosen: escrow across and the pool up, so that the wallet is what
% is left and the funds that add up are a triangle. One stroke for each
% kind of move the walk made.
% Grammar rejected: an equilateral plot of three shares, where a row of
% equal pool is no longer level and the moves out of the triangle have
% nowhere to go; the table itself, which is 211 rows.
% Five-second test: on which side does a stroke cross the dotted line?
\begin{figure}[H]
\centering
\pdfigurehue{pdcobalt}%
\begin{tikzpicture}[pd figure,x=.5cm,y=.5cm,
    row move/.style={draw=pdinkmuted!75,line width=.5pt},
    rising move/.style={draw=pdfocus,line width=.35pt,-{Stealth[length=2.2pt,width=1.8pt]},shorten >=1.6pt},
    stray move/.style={draw=pderror,line width=.7pt,-{Stealth[length=3pt,width=2.6pt]},shorten >=1.8pt},
    place/.style={circle,fill=pdink,inner sep=0pt,minimum size=2.6pt},
    beyond/.style={circle,draw=pderror,fill=pdpage,line width=.6pt,inner sep=0pt,minimum size=3.2pt}]
  % The tables are the recorded runs, read as they were printed: the first
  % line is the command, and every remark starts with '#'.
  \pgfplotstableread[skip first n=1]{excerpts/states-each.txt}\pdmovesreal
  \pgfplotstableread[skip first n=1]{excerpts/states-each-overdraw.txt}\pdmovesloose
  % Counts taken from the tables as they are drawn.
  \ifdefined\pdmovesin\else \newcount\pdmovesin \newcount\pdmovesout \fi
  % One panel: the moves of a table, then the places on top of them.
  \newcommand{\pdmovespanel}[1]{%
    \global\pdmovesin=0 \global\pdmovesout=0
    \pgfplotstablegetrowsof{#1}
    \pgfmathtruncatemacro{\pdlastmove}{\pgfplotsretval-1}
    % Moves along a row first, so that the moves up lie over them.
    \foreach \pass in {0,1,2}{
      \foreach \i in {0,...,\pdlastmove}{
        \pgfplotstablegetelem{\i}{e0}\of#1 \let\ea\pgfplotsretval
        \pgfplotstablegetelem{\i}{p0}\of#1 \let\pa\pgfplotsretval
        \pgfplotstablegetelem{\i}{e1}\of#1 \let\eb\pgfplotsretval
        \pgfplotstablegetelem{\i}{p1}\of#1 \let\pb\pgfplotsretval
        \pgfplotstablegetelem{\i}{wrong}\of#1 \let\wrong\pgfplotsretval
        \ifnum\wrong=1
          \ifnum\pass=2 \global\advance\pdmovesout 1
            \draw[stray move] (\ea,\pa) -- (\eb,\pb); \node[beyond] at (\eb,\pb) {};\fi
        \else\ifnum\pa=\pb
          \ifnum\pass=0 \global\advance\pdmovesin 1 \draw[row move] (\ea,\pa) -- (\eb,\pb);\fi
        \else
          \ifnum\pass=1 \global\advance\pdmovesin 1 \draw[rising move] (\ea,\pa) -- (\eb,\pb);\fi
        \fi\fi
      }
    }
    % Every place where the funds add up to 6.
    \foreach \p in {0,...,6}{
      \pgfmathtruncatemacro{\most}{6-\p}
      \foreach \e in {0,...,\most}{ \node[place] at (\e,\p) {}; }
    }
    % Past this line a project would hold more than was topped up.
    \draw[pd guide] (6.9,-.4) -- (-.4,6.9);
    \foreach \n in {0,2,4,6}{
      \node[pd label,text=pdinkmuted,anchor=north] at (\n,-.55) {\n};
      \node[pd label,text=pdinkmuted,anchor=east] at (-.55,\n) {\n};
    }
    \node[pd label,anchor=north] at (3,-1.5) {units in escrow};
  }
  \pdmovespanel{\pdmovesreal}
  \node[pd label,rotate=90,anchor=south] at (-1.6,3) {units in the pool};
  % Each panel's words sit above it, clear of the strokes and of the
  % other panel.
  \node[pd title,anchor=south west] at (-.6,8.9) {the kernel's keeper};
  \node[pd label,anchor=north west,align=left] at (-.6,8.9)
    {\the\pdmovesin\ kinds of move,\\none past the line};
  \begin{scope}[xshift=4.4cm]
    \pdmovespanel{\pdmovesloose}
    \node[pd title,anchor=south west] at (-.6,8.9) {with the check of the wallet removed};
    \node[pd label,anchor=north west,align=left] at (-.6,8.9)
      {\the\pdmovesout\ more, each to a place\\with more than 6 units};
  \end{scope}
\end{tikzpicture}
\caption{\textbf{No move the kernel makes takes a project's funds past what
was topped up.} A project has 6 units, and the wallet holds whatever is not
in escrow or in the pool. Grey strokes are bonds posted and refunded, which
move along a row. Arrows are slashes, which move up a row and never down.
On the right the core is built without the guard of
Listing~\ref{alg:bonded-build-covered}, and each red arrow posts a bond the
wallet cannot cover. [internal, \texttt{harbor states -{}-each}]}
\label{fig:bonded-moves}
\end{figure}


For a fixed supply, a project's funds are a point in a triangle. Escrow is
on the horizontal axis and the pool on the vertical axis, and the wallet
holds the rest. Conservation says that an operation never leaves the
triangle. The figure shows a little more. Posting or refunding a bond
moves the point along its row. A slash moves it up, and no line goes down,
because no operation takes units out of the pool.
Exercise~\ref{ex:bonded-build-triangle} counts the points, and asks why
the test visits many more states than there are points.

\paragraph{Any amount.}
The exhaustive test stops at 6 units, and the theorem does not. Its proof
is an induction. A project starts at $(0,0,0)$, and only a top-up changes
the total. The step of that induction is about one operation from any
funds, and a bounded model checker can check it completely.
Kani~\pdcite{aws2022kani} compiles a Rust function and asks a solver
whether any input makes an assertion in it fail.
Listing~\ref{alg:bonded-build-kani-step} is the step, written as a Kani
harness. Its \texttt{any\_bonds} is funds and a bond with every field
chosen by the solver. One assumption is made about them: the funds add up
to no more than the ceiling, $2^{53}-1$. \texttt{any\_operation} is one
of the four operations, with any amounts. A \texttt{cover} is the opposite
of an assertion. It requires that some input reaches its line, so the
harness cannot pass just because it had no cases to check.

\pdlistingroom{21}
\pdexcerpt[caption={The step of the induction, as a Kani
harness.},label={alg:bonded-build-kani-step}]{ch07-kani-step}

Kani's results for it and for two more harnesses are in
Listing~\ref{alg:bonded-build-kani}.

\pdlistingroom{9}
\pdrunlines[caption={Kani~0.68 on the core's rules for
funds.},label={alg:bonded-build-kani}]{kani-bonds}{1-9}{1\textendash 9}

Every number in these harnesses is 64 bits that the solver chooses, so for
one operation the bound of the exhaustive test is removed. With the base
case, that a project with no funds holds no units, the result for
\texttt{every\_operation\_conserves\_units} gives
Theorem~\ref{thm:conservation} for this code and any number of operations.
\texttt{two\_bonds\_cannot\_overdraw\_a\_wallet} is
Theorem~\ref{lem:no-overdraft} for two posts made one after the other. The
second result is a control. Its harness makes the step's assertion about
a post that skips the check of the wallet. That assertion is false, and
Kani reports the failure.

\paragraph{Each guard removed.}
The core has five guards on funds, and a second build for each of them
with that guard removed. The same harnesses run on every build.

\pdlistingroom{15}
\pdrunlines[caption={The same harnesses, on five builds of the core with
one guard removed from each.},label={alg:bonded-build-kani-guards}]{kani-bonds}{1,10-23}{1 and 10\textendash 23}

Every build fails at least one harness, so every guard is needed. The
build without the check of the wallet also fails the run with two workers
(Listing~\ref{alg:bonded-build-posts-overdraw}).

\pdlistingroom{5}
\pdrunlines[caption={Two workers and one wallet, with the check of the
wallet removed.},label={alg:bonded-build-posts-overdraw}]{posts-overdraw}{1,6-9}{1 and 6\textendash 9}

The kernel's row is now the same as the row above it. That build is also
the right-hand panel of Figure~\ref{fig:bonded-moves}.

One entry of Listing~\ref{alg:bonded-build-kani-guards} needs an
explanation. Under \texttt{settle-twice}, the build that settles a bond a
second time, the step of the induction is not among the failures. The
next two runs show why. Wren posts a bond of 3 and Heron one of 6, and
Wren requests a refund of hers twice. The kernel as built refuses the
second request.

\pdlistingroom{16}
\pdrun[caption={A refund requested twice. The second request is refused,
and the other bond is refunded in
full.},label={alg:bonded-build-refund-once}]{refund-twice-core}

The same run on the build without the guard:

\pdlistingroom{17}
\pdrun[caption={The same run with the guard on a bond's status removed.
Units are still
conserved.},label={alg:bonded-build-refund-twice}]{refund-twice-settle-twice}

Wren's bond is refunded twice. The second refund is paid from the units
that escrow holds for Heron's bond. The buckets still add up to 10, and
then Heron's refund is refused, because escrow is 3 units short of his
bond. No unit was created or lost, and 3 units that escrow should hold for
Heron's bond are in the wallet. Conservation is a statement about the
total. It does not say which bucket each unit is in. That a bond is
settled once is a second property, and it needed its own harness, which is
the one that fails. \S\ref{sec:conservation} says the same of positive
escrow and conservation: a ledger can have either without the other.
Exercise~\ref{ex:bonded-build-twice} asks what is false of the funds and
the bonds together after step 5.
\pdscope{The exhaustive test covers amounts of 1 to 3, with at most 6
units and 3 bonds, and it compares the ledger with a reference model
written by the same author as the core, from the same proof sketch. The
harnesses cover one operation from any funds, or two in a row, on the four
functions that \texttt{decide} calls. They do not cover the TypeScript
that stores each result.}

\begin{pdboundary}
\setlength{\parskip}{4pt plus 1pt}
These funds are in memory, and they are lost when the process stops.
Chapter~1's kernel stores its claims in a database. This layer does not
use it yet.

Callers are not authenticated. Any caller can top a project up, and can
refund or slash any bond. Definition~\ref{def:settlement} requires an
authorized ruling before a settlement, and this build has no way to give
one.

A bond is not linked to anything. It is not posted for a claim or for a
piece of work, and no rule requires a worker to hold one before it starts.
A worker can still claim a path at no cost. The next section links a bond
to a claim.

Nothing is paid out of the pool. A forfeit goes there and stays, so no
pair of parties can get it back
(Property~\ref{prop:third-party-forfeiture}). Who the pool is for is not
decided here.

A unit is a number in the kernel's state. Whether real money backs it is
a question of custody, which this code does not address.
\end{pdboundary}
